\section{Position embeddings 与 RoPE}

本节解释从绝对位置到 RoPE 的迁移：相对位置信息通过 Q/K 旋转进入 attention score。


\begin{importantbox}{RoPE 的一句话定义}
RoPE 把位置信息编码为 query/key 向量的旋转，使 attention score 中自然出现相对位置。它不是把 position embedding 加到 token embedding 上，而是在每层 attention 的 Q/K 空间里改变几何关系。
\end{importantbox}

\begin{knowledgebox}{RoPE 公式链路}
把每两个 hidden dimensions 看成一个二维平面。位置 \(i\) 上的向量对 \((x_{2k},x_{2k+1})\) 被旋转角度 \(i\theta_k\)：
\[
\begin{bmatrix}
x'_{2k}\\
x'_{2k+1}
\end{bmatrix}
=
\begin{bmatrix}
\cos(i\theta_k) & -\sin(i\theta_k)\\
\sin(i\theta_k) & \cos(i\theta_k)
\end{bmatrix}
\begin{bmatrix}
x_{2k}\\
x_{2k+1}
\end{bmatrix}.
\]
当 query 和 key 都按各自位置旋转后，它们的点积会依赖相对位移 \(i-j\)。这就是 RoPE 能支持 relative-position behavior 的几何原因。
\end{knowledgebox}

\begin{figure}[H]
\centering
\includegraphics[width=0.88\textwidth]{slide-029.jpg}
\caption{Slide 30: Position embedding variants}
\figsource{CS336 Spring 2026 Lecture 3 官方幻灯片。}
\end{figure}

\noindent\textbf{展开说明：}本页列出 sinusoidal、learned absolute、relative、ALiBi、RoPE 等位置编码路线。


\begin{figure}[H]
\centering
\includegraphics[width=0.88\textwidth]{slide-030.jpg}
\caption{Slide 31: RoPE high-level}
\figsource{CS336 Spring 2026 Lecture 3 官方幻灯片。}
\end{figure}

\noindent\textbf{展开说明：}本页给出 RoPE 的相对位置目标：让 attention score 能感知相对距离。



“希望 attention 依赖相对位置”还不是构造。Slide 32 给出 RoPE 的几何桥梁：内积对共同旋转不变，因此可以把 token 在不同绝对位置上的 Q/K 看成同一内容向量经过不同角度旋转。比较两个 token 时，绝对角度相消，只留下旋转角之差，也就是位置差。图中的短语例子是在说明同一个词移到不同上下文位置后，内容坐标和位置坐标如何解耦。

\begin{figure}[H]
\centering
\includegraphics[width=0.88\textwidth]{slide-031.jpg}
\caption{Slide 32：利用旋转后的内积把绝对位置转化为相对位置差。}
\figsource{CS336 Spring 2026 Lecture 3 官方幻灯片。}
\end{figure}

\begin{importantbox}{从共同旋转到相对位移}
若位置 $i$ 和 $j$ 的向量分别乘旋转矩阵 $R_i$ 与 $R_j$，则
\[
(R_i q)^\top(R_j k)=q^\top R_i^\top R_j k=q^\top R_{j-i}k.
\]
最后一项只依赖 $j-i$。这条等式解释了 RoPE 为什么把位置注入 Q/K 而不是 V，也解释了 attention score 如何自然获得方向和距离信息。
\end{importantbox}

高维向量可以用许多旋转方式，Slide 33 说明标准 RoPE 的具体选择：把坐标两两配对，每一对在二维平面中旋转，并给不同维度对分配不同频率。高频平面敏感于局部位移，低频平面覆盖更长距离；这种多尺度结构类似 sinusoidal embeddings，但作用位置从输入加法变为 attention 的 Q/K 几何。

\begin{figure}[H]
\centering
\includegraphics[width=0.88\textwidth]{slide-032.jpg}
\caption{Slide 33：标准 RoPE 将坐标两两配对，并在多个二维平面中使用不同旋转频率。}
\figsource{CS336 Spring 2026 Lecture 3 官方幻灯片；Su et al. 2021。}
\end{figure}

\begin{warningbox}{RoPE 公式成立，不代表长度外推自动成立}
相对位置形式保证了几何结构，却没有保证训练长度之外的频率仍被模型正确使用。改变 base、缩放位置或只旋转部分维度都会改变外推行为；长上下文 recipe 必须在目标长度上测 token-level cross-entropy、retrieval 和 attention pattern，而不是只检查公式可运行。
\end{warningbox}


\begin{figure}[H]
\centering
\includegraphics[width=0.88\textwidth]{slide-033.jpg}
\caption{Slide 34: Actual RoPE math}
\figsource{CS336 Spring 2026 Lecture 3 官方幻灯片。}
\end{figure}

\noindent\textbf{展开说明：}本页给出 RoPE 公式：把偶/奇维组成二维平面，用 sin/cos 对每对维度旋转。


\begin{figure}[H]
\centering
\includegraphics[width=0.88\textwidth]{slide-034.jpg}
\caption{Slide 35: RoPE implementation}
\figsource{CS336 Spring 2026 Lecture 3 官方幻灯片。}
\end{figure}

\noindent\textbf{展开说明：}本页说明 RoPE 在 attention 的 Q/K 上实施，通常不直接改 V。实现细节会影响 shape 和广播。

\subsection{本章小结}
本节的关键不是记住所有模型名，而是把公开模型的设计选择转化成可迁移判断：共识默认值可以作为起点，例外需要知道原因，涉及系统和推理成本的选择必须结合资源账本讨论。

